\documentclass[10pt,a4paper]{amsart}
\usepackage[margin=19mm]{geometry}
\usepackage{amsmath,amssymb,mathtools}
\usepackage[scr=rsfs,cal=euler]{mathalfa}
\usepackage{xfrac}
\usepackage[unicode,hidelinks,bookmarks=false]{hyperref}
\newcommand{\ie}{i.e.\ }
\newcommand{\cf}{cf.\ }
\newcommand{\resp}{respectively\ }
\newcommand{\Subsection}{\subsection}
\providecommand{\nobreakdash}{\nobreak\hbox{-}}
\setlength{\emergencystretch}{2em}
\multlinegap0pt
\usepackage{amssymb}
\usepackage{mathtools}
\usepackage[shortlabels]{enumitem}
\newtheorem{theorem}{Theorem}
\newcommand{\ao}{\mathrm{ao}}
\newcommand{\hK}{\mathbb K}
\newcommand{\tri}{\mathbin{\triangleleft}}
\title{The \(V_n\) Invariants as Colored Links--Gould Invariants:\\ A Root--Center Approach}
\author{Jiuhe Liu}
\date{Source: arXiv:2608.28720v2 (CC BY 4.0)}
\begin{document}
\maketitle

\noindent\emph{Source and licence.} The \(V_n\) Invariants as Colored Links--Gould Invariants:\\ A Root--Center Approach. Version arXiv:2608.28720v2 (CC BY 4.0), distributed by arXiv under the Creative Commons Attribution 4.0 International licence, http://creativecommons.org/licenses/by/4.0/, as recorded in the arXiv licence metadata for this exact version and shown on the abstract page https://arxiv.org/abs/2608.28720v2. Full licence notice: \texttt{licenses/arxiv-2608.28720-CC-BY-4.0.txt}.\par\smallskip

\noindent\textit{Editorial scope.} Counted unit: the article's theorem thm:object (source0.tex lines 870--878). The article's complete original proof of it is delivered with it (source0.tex lines 880--931, one \texttt{proof} environment of 52 lines). No mathematics is written, repaired or re-derived: the statement and the proof are the article's own lines, in the article's own notation, and no retained line is substituted or re-flowed.  Retained uncounted as the source-local context the proof needs: lines 355--361; lines 847--854; lines 857--861.  Nothing was omitted from any retained span, so the package carries no omission record.  No retained line contains a citation, so the wrapper carries no bibliography and no bibliography entry.  No retained line contains a figure environment, an image-inclusion command or drawing code, so no image is delivered and the package declares no asset dependency.  Editorial formatting only, in addition: an AMS \texttt{amsart} preamble that restates the article's own declarations and macros used by the retained lines, namely the environments \texttt{theorem} on separate counters, together with the standard packages those lines need.  The retained environments print in this document as Theorem 1 in order of appearance, whereas the article prints the same items with the numbering of its own full document.  The complete native source of the article as distributed in the arXiv e-print for this exact version is bundled unchanged as \texttt{sources/208804/source0.tex}.
\begin{equation}\label{eq:Yn-basis}
\begin{split}
 \{1,v_n\}
 &\cup\{(x_1x_2)^k,(x_2x_1)^k:1\leq k<n\}\\
 &\cup\{(x_1x_2)^kx_1,(x_2x_1)^kx_2:0\leq k<n\}.
\end{split}
\end{equation}

\begin{equation}\label{eq:root-center-actions}
\begin{aligned}
 y\tri\mathsf h_1&=(n(\beta+1)-b)y,&
 y\tri\mathsf h_2&=-(n\beta+a)y,\\
 y\tri\mathsf c_1&=-2n(\beta+1)y,&
 y\tri\mathsf c_2&=2n\beta y.
\end{aligned}
\end{equation}

\begin{equation}\label{eq:highest-weight}
 w_{n,\beta}\tri E_1=w_{n,\beta}\tri E_2=0,\qquad
 w_{n,\beta}\tri\mathsf h_1=n\beta w_{n,\beta},\qquad
 w_{n,\beta}\tri\mathsf h_2=-n(\beta+1)w_{n,\beta}.
\end{equation}

\begin{theorem}[Object identification]\label{thm:object}
For \(\beta\neq0,-1\),
\begin{equation}\label{eq:object}
 \mathcal Y_n
 \cong W_{n,\beta}^{\ao}\boxtimes L_{n,\beta}
\end{equation}
as a right module over the twisted paired double.  The root factor
\(W_{n,\beta}^{\ao}\) is simple and has dimension \(4n\).
\end{theorem}

\begin{proof}
Write
\[
 A_k=(x_1x_2)^k,\quad B_k=(x_2x_1)^k,\quad
 C_k=(x_1x_2)^kx_1,\quad D_k=(x_2x_1)^kx_2.
\]
The vector \(v_n=A_n+tB_n\) is even, is killed by \(E_1,E_2\), and has
the weight in \eqref{eq:highest-weight}.

The negative root actions differ by nonzero degree scalars from braided
derivatives \(\partial_1,\partial_2\).  They obey
\[
\begin{array}{ll}
\partial_1(A_k)=-Q^kD_{k-1},&
\partial_2(A_k)=C_{k-1},\\
\partial_1(B_k)=D_{k-1},&
\partial_2(B_k)=-Q^kC_{k-1},
\end{array}
\]
and
\[
 \partial_1(v_n)=(t-Q^n)D_{n-1},\qquad
 \partial_2(v_n)=(1-tQ^n)C_{n-1}.
\]
Both coefficients are nonzero precisely under
\(\beta\neq0,-1\).  Furthermore,
\[
 \partial_1(C_k)=A_k+Q^kB_k,\qquad
 \partial_2(D_k)=Q^kA_k+B_k.
\]
The coefficient matrix has determinant \(1-Q^{2k}\neq0\).
Descending induction therefore produces every basis vector in
\eqref{eq:Yn-basis} from \(v_n\); hence \(v_n\) is cyclic.

It remains to show that \(v_n\) is the unique highest vector.  On
\(\operatorname{span}\{A_k,B_k\}\), \(1\leq k<n\), the common kernel of
\(E_1,E_2\) is controlled by
\[
 \begin{pmatrix}
 1&-t_{1,n}Q^k\\
 -t_{2,n}Q^k&1
 \end{pmatrix},
\]
whose determinant is \(1-Q^{2(k-n)}\neq0\).  The corresponding odd-degree
formulas show the same for \(\operatorname{span}\{C_k,D_k\}\), while the
top common kernel is exactly \(\hK v_n\).  Any nonzero submodule contains
a vector of maximal degree, hence a highest vector, and therefore
contains \(v_n\).  Cyclicity then forces the submodule to be all of
\(\mathcal Y_n\).  This proves simplicity and identifies the root factor.
The central action in \eqref{eq:root-center-actions} gives the tensor
factor \(L_{n,\beta}\), proving \eqref{eq:object}.
\end{proof}

\end{document}
